% zakljucak.

\section{Zaključak}
\label{sec:zakljucak}

U radu su riješena četiri zadatka koja pokrivaju raspon od uvjeta optimalnosti
za kvadratni program do sinteze regulatora za hibridni sustav.

U Zadatku~1 pokazano je da je zadani problem strogo konveksan kvadratni program,
pa su KKT uvjeti dovoljni za globalnu optimalnost. Analitičko rješenje
$x^\star = (190/33,\, 140/33)$ uz $\lambda^\star = (404/33,\, 0)$ potvrđeno je
trima neovisnim numeričkim putevima --- funkcijom \texttt{quadprog}, YALMIP-om i
rješavanjem dualnog problema --- uz dualni jaz reda $10^{-8}$. Osjetljivost
optimalne vrijednosti numerički se poklopila s $-\lambda^\star$, čime je
potvrđena interpretacija množitelja kao "cijene" ograničenja.

U Zadatku~2 pokazano je da intervalna nesigurnost koeficijenata ne mijenja klasu
problema: najgori slučaj po kvadru svodi se na član $r^{\mathsf T}|x|$, koji je
konveksan, pa robusna protuformulacija ostaje kvadratni program. Rezultat je
poučan --- nominalno rješenje, koje leži točno na granici, nedopustivo je u
$49{,}9\,\%$ slučajnih realizacija, dok robusno nije prekršilo ograničenja ni
jednom. Cijena te sigurnosti je porast funkcije cilja za oko $40\,\%$.

Zadatak~3 pokazao je koliko formulacija određuje ishod. Budući da je trošak
proporcionalan količini materijala, a ne njezinu kvadratu, prirodna mjera je
$\ell_1$ norma, što nakon epigrafske reformulacije daje linearni program.
Usporedba slučajeva pokazala je da strože ograničenje na promjenu nagiba
poskupljuje izvedbu za oko $28\,\%$ i mijenja strukturu rješenja: umjesto
ograničenja nagiba, rješenje počinje oblikovati ograničenje zakrivljenosti.
Prijelaz s $N = 20$ na $N = 100$ promijenio je trošak za tek oko $2\,\%$, pa je
gruba diskretizacija dovoljna za inženjersku procjenu.

Zadatak~4 dao je najizrazitiji rezultat: oba dinamička moda imaju identične
svojstvene vrijednosti $-1 \pm \mathrm{i}\pi/2$ i oba su Hurwitzova, a prekidački
sustav je ipak nestabilan, sa spektralnim radijusom monodromijske matrice
$1{,}218$. Time je pokazano da se stabilnost hibridnih sustava ne može ispitivati
mod po mod. Statički regulator $K = [-20,\; -4\pi/3]$ konstruiran je tako da uz
$P = I$ zadovoljava Ljapunovljevu nejednakost za oba moda, čime je zajamčena
eksponencijalna stabilnost \emph{za proizvoljno prekapčanje}, a ne samo za
zadani periodični raspored.

Zajednički zaključak jest da konveksnost nije tehnička sitnica nego glavni alat.
Ona je omogućila da se za Zadatak~1 tvrdi globalna optimalnost, da robusna
protuformulacija u Zadatku~2 ostane jednako lako rješiva kao nominalni problem,
da se Zadatak~3 svede na linearni program, i da se u Zadatku~4 problem sinteze
regulatora --- inače bilinearan --- supstitucijom $Q = P^{-1}$, $Y = KQ$ prevede
u konveksan skup linearnih matričnih nejednakosti.
